\chapter{Lo mínimo de Lean: cuatro palabras}
\label{cap:lean}

Este libro está lleno de código Lean 4 y no es un libro sobre Lean. Para leerlo hacen falta cuatro
palabras. Con ellas se entiende cada fragmento que aparece; sin ellas, ninguno. Todo lo demás
—tácticas, elaboración, instancias— es maquinaria que el lector puede ignorar sin perder el hilo.

\section{\ident{inductive}: cómo se dice «esto es todo lo que hay»}

Un tipo inductivo se declara enumerando sus \emph{constructores}, y esa enumeración es exhaustiva:
no hay más términos que los que se obtienen aplicándolos.

\leanfrag{fol_Term}

\noindent Un término es una variable o la aplicación de un símbolo de función a una lista de
términos. Nada más. De esa declaración salen gratis tres cosas que en la teoría objeto cuestan
axiomas:

\begin{center}\small
\begin{tabularx}{\linewidth}{@{}l >{\raggedright\arraybackslash}X@{}}
\toprule
constructores disjuntos & una variable nunca es una aplicación \\
constructores inyectivos & si dos aplicaciones son iguales, coinciden símbolo y argumentos \\
inducción estructural & para demostrar algo de todo término basta hacerlo caso por constructor \\
\bottomrule
\end{tabularx}
\end{center}

\begin{leccion}[Lo que regala el compilador, en primer orden se paga]
Esas tres cosas son exactamente lo que la parte~\ref{parte:aritmetizacion} tiene que \emph{demostrar} cuando reconstruye
listas y árboles sintácticos dentro de la aritmética, donde no hay tipos inductivos: hay números.
Ver hecho a mano lo que aquí es una línea enseña qué había realmente dentro de esa comodidad.
\end{leccion}

\section{Las tres palabras que son tres categorías morales}

Las otras tres —\ident{def}, \ident{theorem}, \ident{axiom}— parecen tres formas de escribir lo
mismo. No lo son. Son tres compromisos distintos, y la diferencia entre ellos organiza el proyecto
entero.

\subsection*{\ident{def} — nombrar no cuesta nada}

Una definición introduce una abreviatura. No afirma nada y no puede ser falsa: como mucho, puede ser
inútil.

\leanfrag{consN}

\subsection*{\ident{theorem} — se gana}

Un teorema afirma algo, y sólo entra en el sistema si el núcleo de Lean acepta su demostración.

\leanfrag{two_mul_consN}

\subsection*{\ident{axiom} — es una deuda}

Un axioma afirma algo \textbf{sin demostración}. Lean lo acepta y lo apunta. Es legítimo cuando lo
que se postula es una regla del juego o un principio que la teoría de destino tiene por definición
— y es peligroso siempre, porque un axioma falso hace demostrable cualquier cosa.

\leanfrag{fol_gen}

\begin{muro}[Un axioma falso ya ha pasado aquí, dos veces]
En \modulo{FOL} hubo un \ident{axiom} cuyo enunciado era literalmente falso, con contraejemplo de
tres líneas; hoy es un teorema con el enunciado corregido. Y en la teoría objeto hubo otro que hacía
la teoría inconsistente durante meses. Los dos episodios están contados en este libro, y los dos
empiezan igual: alguien escribió \ident{axiom} donde debía haber escrito una pregunta.
\end{muro}

\section{Por eso hay una línea gris debajo de cada teorema}

Si un axioma es una deuda, lo mínimo es llevar la cuenta. Lean lo permite: dado cualquier teorema,
se puede preguntar de qué axiomas depende — no los que están declarados en el fichero, sino los que
su demostración usa realmente, siguiendo toda la cadena. Esa lista es el \defterm{footprint}, y en
este libro se imprime debajo de cada enunciado citado:

\footprint{\fpi{propext}, \fpi{Classical.choice}, \fpi{Quot.sound}}

\begin{leccion}[Cómo se lee]
Los tres de arriba son la base del propio Lean y salen prácticamente en todo. Lo interesante es lo
que aparece \emph{además}: si un teorema cita un esquema de inducción, es que lo necesita; si no lo
cita, es que no. Cuando el footprint no añade nada a la base ya aceptada se dice que el resultado es
\defterm{net-0}. Y cuando aparece algo que no debería, hay un problema — así se descubrió la
inconsistencia de la parte~\ref{parte:noescrito}.
\end{leccion}

\section{Dos palabras más, que no son de Lean sino de este proyecto}

Aparecen a menudo y conviene fijarlas aquí. Un \defterm{sondeo} es un experimento compilado que
contesta una pregunta —muchas veces respondiendo \emph{que no}— y que vive fuera de la compilación
principal. La \defterm{cuarentena} es el directorio donde se apartaron los módulos que dependían de
la inconsistencia; hoy está vacío, y el episodio es la Parte~\ref{parte:noescrito}.

\begin{leccion}[Todo lo anterior en una línea]
\ident{def} nombra, \ident{theorem} se gana, \ident{axiom} se debe. El resto de este libro consiste,
básicamente, en pagar deudas y en anotar las que no se han podido pagar todavía.
\end{leccion}
